% ---- BEGIN monograph/frontmatter/abstract.tex ----
\chapter*{Abstract}
\addcontentsline{toc}{chapter}{Abstract}

This work develops a conditional research programme in which temporal structure is constructed from an explicitly declared relational model before spacetime geometry is introduced. The starting data are relational states and directed transitions carrying a normalized probability distribution, a normalized complex representative on phase-bearing sectors, positive transition weights, an ordering parameter distinct from metric clock time, and an inherited dimensionless normalization. The programme then asks which temporal quantities follow from those inputs by definition, exact identity, dynamical closure, or calibrated physical bridge.

Directed transitions carry both an exact Shannon entropy difference and a geometric phase link. Their combination with a separately typed non-exact path-affinity term is used to construct an auditable temporal transition phase, while closed-cycle identities separate exact state-function changes from genuinely non-exact directional information. Positive transition activity then generates an elapsed measure and a relative lapse before a spacetime metric is introduced. Downstream layers develop NOW support, bifurcation, transport, memory, constructive retrodiction, preregistered physical tests, and finally a temporal-to-spacetime closure interface.

The use of the word \emph{derived} is deliberately restricted: IDT does not claim that physical time follows uniquely from Shannon information alone. Its claims are conditional on the declared relational primitives and domains. Standard results from information theory, geometric phase, nonequilibrium response, quantum-state geometry, general relativity, and horizon thermodynamics are treated as imported background or recovery benchmarks rather than as new predictions.

Formal derivation, computational evidence, and physical claim status are maintained as separate layers. Exact identities, reference-model reconstructions, calibrated physical correspondences, and replicated measured results are not interchangeable evidence classes. The strongest present results are constructive and algebraic within their stated domains; the strongest physical claims remain attached to explicit falsification protocols whose calibration variables and null models are fixed independently of the measured outcome.

% ---- END monograph/frontmatter/abstract.tex ----

% ---- BEGIN monograph/frontmatter/methodology.tex ----
\chapter*{Methodological Note}
\addcontentsline{toc}{chapter}{Methodological Note}

The dependency rule of this programme is
\[
\boxed{\text{derive temporal structure first; introduce spacetime later}.}
\]
The construction first imports the TIR relational root
\[
\mathfrak Z_{\rm rel}
\to\mathcal P
\to\{N,S\}
\to\frac12
\to\ln2
\to\mathbb C^2,
\]
without identifying relational zero with a time coordinate, vacuum state, or initial event. IDT-local construction begins only after that upstream packet is admitted, from relational information and phase transport. Each new temporal object is admitted only after its upstream dependencies are explicit. The intended order is
\[
\begin{aligned}
\text{TIR relational root}
&\to\text{information}
\to\text{directed transition}
\to\text{phase structure}
\to\text{temporal state}\\
&\to\text{NOW/bifurcation}
\to\text{transport/memory}
\to\text{metric calibration}\\
&\to\text{space}
\to\text{Einstein closure}.
\end{aligned}
\]

Mathematical construction, computational evidence, and physical claim status are kept distinct. A candidate is first assigned typed variables, equations, dependencies, and falsification conditions. Algebraic identities are proved where possible; numerical controls are used where analytic closure is not yet available. Only then can a claim move to a stronger registry status.

The ordering parameter used in the early formalism is kept distinct from metric clock time. Metric calibration is therefore a downstream bridge rather than a hidden starting assumption. Likewise, chirality, topology, fractal scaling, memory, and subjective temporal accumulation are treated as derived-property targets with their own observables and tests.

\section*{Prior-art and novelty discipline}

Whenever the construction invokes an established object, the corresponding chapter cites the standard source or a primary historical reference. Shannon entropy, Kullback--Leibler divergence, Pancharatnam/Berry phase, Onsager response, reversible Markov entropy gradients, Noether conservation, Einstein gravity, Bell/CHSH tests, and Hawking/Gibbons--Hawking horizon thermodynamics are therefore treated as external scientific context rather than as IDT discoveries.

A statement is called \emph{IDT-specific} only when it depends on the programme's own typed composition, decoder, activity measure, memory carrier, calibration map, or falsification rule. Recovery of a known physical equation is labelled a recovery benchmark. A physical discriminator requires an independently determined IDT quantity that can be evaluated before the target measurement is inspected.

\section*{Falsification before promotion}

For a proposed physical correspondence, the calibration map and null model must be fixed independently of the target outcome. If the same measured quantity is used both to define the IDT predictor and to confirm it, the exercise is classified as calibration or consistency checking rather than discrimination. Promotion to a physical result requires immutable raw data, uncertainty propagation, declared exclusions, and independent replication where the claim warrants it.
% ---- END monograph/frontmatter/methodology.tex ----

% ---- BEGIN monograph/frontmatter/research_positioning.tex ----
\chapter*{Scope, Primitives, and Claim Semantics}
\addcontentsline{toc}{chapter}{Scope, Primitives, and Claim Semantics}

\section*{What is and is not being derived}

Throughout this monograph, the word \emph{derived} has a conditional mathematical meaning. Temporal structures are derived from a declared relational model and its typed primitives; they are not claimed to follow uniquely from Shannon information alone, nor is the existence of physical time inferred from an information measure without additional assumptions.

Before any IDT-local temporal object is introduced, the branch imports the typed TIR relational root
\[
\boxed{
\mathfrak Z_{\rm rel}=(\varnothing,\varnothing)
\to\mathcal P
\to\{N,S\}
\to\frac12
\to\ln2
\to\mathbb C^2.
}
\]
Here \(\mathfrak Z_{\rm rel}\) is a pre-object dependency boundary, not \(t=0\), not a vacuum state, and not an initial event. Conditional on that imported root, the IDT-local temporal branch admits the following objects:
\begin{enumerate}
\item a set of relational states and admitted directed transitions;
\item a normalized probability vector \(p(s)\) on each state, with Shannon entropy \(H_S\);
\item a normalized complex representative \(\psi(s)\) on sectors where phase transport is defined;
\item positive admitted transition probabilities or rates, together with declared relational mobility fields;
\item an ordering parameter \(\lambda\), which is \emph{not} metric clock time;
\item the inherited normalization \(\kappa=\ln 2/(24\pi)\).
\end{enumerate}

From those inputs the programme asks which further structures follow by definition, algebraic identity, dynamical closure, or experimentally testable bridge. The positive elapsed measure, relative lapse, memory carriers, retrodiction maps, and later spacetime interface are therefore downstream constructions whose validity is conditional on their stated domains and dependencies.

\section*{Standard ingredients and the IDT-specific composition}

Several ingredients used in the construction are standard and are not presented as new results. These include Shannon entropy \cite{Shannon1948}, Kullback--Leibler divergence \cite{KullbackLeibler1951}, Pancharatnam/Berry geometric phase \cite{Pancharatnam1956,Berry1984}, quantum-state information geometry \cite{ProvostVallee1980}, Onsager reciprocity and nonequilibrium network structure \cite{Onsager1931,Schnakenberg1976}, and entropy-gradient formulations for reversible finite Markov chains \cite{Maas2011}. The relativistic closure layer likewise uses standard Noether and Einstein structures \cite{Noether1971,Einstein1916}.

The programme-specific content is the typed way these ingredients are composed and constrained: the separation of exact state entropy from non-exact path affinity, the positive-activity temporal measure, the clock-ratio and memory interfaces, the constructive retrodiction carrier, the evidence-class firewall between retrodiction and retrocausal claims, and the explicit temporal-to-spacetime closure gates. A standard identity remains standard even when it is used as an important recovery target inside this architecture.

\section*{Evidence classes}

Every substantive statement belongs to one of five evidence classes:
\[
\boxed{
\begin{array}{c}
\text{definition / admitted primitive}\\
\downarrow\\
\text{proved identity on a declared domain}\\
\downarrow\\
\text{reference-model or software result}\\
\downarrow\\
\text{calibrated physical correspondence}\\
\downarrow\\
\text{replicated measured physical result}.
\end{array}}
\]
A result does not move upward merely because more unit tests pass. Hosted tests establish implementation consistency and regression protection; they are not independent physical observations.

\section*{Theory map}

The monograph follows four layers:
\[
\boxed{
\begin{aligned}
\textbf{I. Foundations:}\quad&
\text{TIR relational root}
\to
\text{relational state}
\to
\text{information + phase}
\to
\text{positive temporal activity},
\\[1mm]
\textbf{II. Temporal dynamics:}\quad&
\text{elapsed measure}
\to
\text{NOW/bifurcation}
\to
\text{transport + memory},
\\[1mm]
\textbf{III. Inference and experiment:}\quad&
\text{retrodiction}
\to
\text{withheld-variable tests}
\to
\text{preregistered physical tests},
\\[1mm]
\textbf{IV. Spacetime closure:}\quad&
\text{temporal coframe}
\oplus
\text{independent spatial/source branch}
\to
\text{Einstein interface}.
\end{aligned}}
\]

This map also fixes the novelty boundary. Standard-physics relations used at Layer IV are recovery benchmarks unless the temporal branch supplies an independently determined quantity that can disagree with the standard prediction before the measurement is inspected.
% ---- END monograph/frontmatter/research_positioning.tex ----

% ---- BEGIN monograph/frontmatter/symbols.tex ----
\chapter*{Notation and Principal Symbols}
\addcontentsline{toc}{chapter}{Notation and Principal Symbols}

\small
\setlength{\LTpre}{0pt}
\setlength{\LTpost}{0pt}
\begin{longtable}{p{0.20\textwidth}p{0.27\textwidth}p{0.43\textwidth}}
\toprule
\textbf{Symbol} & \textbf{Type / domain} & \textbf{Role} \\
\midrule
\endfirsthead
\toprule
\textbf{Symbol} & \textbf{Type / domain} & \textbf{Role} \\
\midrule
\endhead
\(s,a,b\) & relational states & State labels; \(a\to b\) denotes an admitted directed transition. \\
\(p(s)\) & probability simplex & Normalized state probability vector. \\
\(H_S(s)\) & bits & Shannon entropy of \(p(s)\). \\
\(\Delta H_e\) & bits & Exact state-function entropy difference along edge \(e\). \\
\(\psi(s)\) & normalized complex vector & Complex representative used for phase transport. \\
\(G_{a\to b}\) & \(U(1)\) & Normalized Pancharatnam link on a non-orthogonal pair. \\
\(\gamma_B(C)\) & angle & Closed-cycle geometric holonomy. \\
\(\sigma_e\) & bits & Non-exact directional information increment / path affinity. \\
\(\kappa\) & dimensionless & Inherited normalization \(\ln2/(24\pi)\). \\
\(L_e\) & \(U(1)\) & Shannon--phase temporal transition link. \\
\(q_e\) & angle & Argument of \(L_e\); transition phase weight. \\
\(\rho_R,\eta_R\) & positive relational fields & Density- and viscosity-like fields controlling the symmetric mobility sector. \\
\(M_{ab}\) & positive scalar & Symmetric pair mobility. \\
\(W_{a\to b}\) & positive rate & Directed transition rate. \\
\(\mathfrak a\) & positive rate & Symmetric activity \(W_++W_-\); generator of elapsed measure. \\
\(\mathfrak j\) & signed rate & Antisymmetric current \(W_+-W_-\). \\
\(\lambda\) & ordering coordinate & Monotone ordering parameter; not metric clock time. \\
\(\Theta\) & elapsed activity measure & Primitive accumulated temporal measure with \(d\Theta=\mathfrak a\,d\lambda\). \\
\(N_R(x|r)\) & positive ratio & Relative lapse / clock-rate ratio \(\mathfrak a_x/\mathfrak a_r\). \\
\(\Phi(N_R)\) & dimensionless & Relative-information clock potential \(N_R-1-\ln N_R\). \\
\(\widehat\tau\) & calibrated time & Physical elapsed time after reference-clock calibration. \\
\(m\) & complex memory coordinate & Two-component memory carrier used in the reference reconstruction class. \\
\(\Delta M_n\) & operator difference & Projective pre/post event memory imprint. \\
\(\alpha\) & ordered index packet & Active-attractor sequence in the stratified decoder. \\
\(\mathcal W\) & winding packet & Oriented winding lineage retained by the compressed decoder. \\
\(\rho_k\) & positive radius & Radial coordinate relative to the active attractor. \\
\(\Theta_R\) & length-valued one-form & Relativistic temporal coframe input \(N_Rc\,dt\). \\
\(g_{\mu\nu}\) & spacetime metric & Introduced only at the downstream spacetime closure layer. \\
\(T_{\mu\nu}\) & stress-energy tensor & Physical source ledger at the Einstein interface. \\
\(\kappa_E\) & gravitational coupling & Standard \(8\pi G/c^4\); not derived anew by IDT. \\
\(\kappa_H\) & surface gravity scale & Horizon parameter used in the standard Euclidean thermal benchmark. \\
\bottomrule
\end{longtable}
\normalsize

Symbols introduced only inside a specialised decoder, experiment, or source adapter are defined locally in the corresponding chapter. Reuse of a symbol in a different typed domain is avoided where it would change the meaning of an established quantity.

% ---- END monograph/frontmatter/symbols.tex ----

% ---- BEGIN monograph/chapters/02_tir_entrypoint.tex ----
\chapter{TIR Entrypoint}

The temporal branch does not begin from a physical time coordinate. It imports only registry-declared TIR objects and preserves their dependency order.

The upstream relational root is
\[
\boxed{
\mathfrak Z_{\rm rel}
=
(\varnothing,\varnothing)
\longrightarrow
\mathcal P=\{p\}
\longrightarrow
\Pi_1=\{N,S\}
\longrightarrow
\frac12
\longrightarrow
\ln2
\longrightarrow
\mathbb C^2.
}
\]

Here $\mathfrak Z_{\rm rel}$ is the TIR empty relational presentation. It is not a point and is not a temporal coordinate. In particular,
\[
\boxed{
0_{\rm rel}
\neq
t=0
\neq
\tau_{\rm int}=0
\neq
\text{vacuum state}
\neq
\text{initial event}.
}
\]
The point is the minimal nonempty support downstream of relational zero. The first binary distinction is the minimal nontrivial distinction; exchange symmetry then fixes the normalized share $1/2$, and Shannon information gives $\ln2$.

IDT imports this root rather than re-deriving or reinterpreting it. The typed cross-repository contract is maintained in
\path{formalism/00A_tir_relational_zero_entrypoint.md}.

The canonical TIR normalization used downstream in IDT is
\[
\kappa=\frac{\ln2}{24\pi}.
\]
Its denominator and flavour/mixing ownership remain TIR-side dependencies; IDT does not promote the normalization from the zero layer alone.

The first local temporal construction combines three typed objects rather than collapsing them into one: Shannon state entropy $H_S$, geometric phase transport $G_e$, and a transition-associated non-exact information-production increment $\sigma_e$.

The dependency chain at this stage is therefore
\[
\boxed{
\text{TIR relational root}
\longrightarrow
\text{relational state}
\longrightarrow
(H_S,\psi)
\longrightarrow
(\Delta H_e,G_e)
\longrightarrow
L_e
\longrightarrow
q_e
\longrightarrow
\mathcal K_T.
}
\]

The first explicit candidate closure of the edge quantity is
\[
\sigma_{a\to b}=\log_2\frac{P(b\mid a)}{P(a\mid b)}.
\]
This closes the free edge parameter at the information-theoretic level while preserving the distinction between state entropy and path affinity. The next upstream debt is to derive the forward/reverse transition weights from the relational density and viscosity fields rather than supplying them externally.

The publication firewall is:
\[
\boxed{
\text{relational zero}
\not\Rightarrow
\text{physical time};
\qquad
\text{TIR structural import}
\not\Rightarrow
\text{IDT empirical promotion}.
}
\]

% ---- END monograph/chapters/02_tir_entrypoint.tex ----

% ---- BEGIN monograph/chapters/01_information.tex ----
\chapter{Information and Relational Order}

The temporal programme begins from relational states rather than from a pre-existing time coordinate. Each state $s$ is assigned a probability distribution
\[
p(s)=(p_1,\ldots,p_m),
\qquad
p_i\geq0,
\qquad
\sum_i p_i=1,
\]
and its Shannon entropy in bits \cite{Shannon1948},
\[
\boxed{
H_S(s)=-\sum_i p_i(s)\log_2p_i(s).
}
\]
The inherited normalization is
\[
\kappa=\frac{\ln2}{24\pi}.
\]

For an admitted directed transition $e:a\to b$, the first informational quantity is the exact state-function difference
\[
\Delta H_e=H_S(b)-H_S(a).
\]
This quantity measures the change in the Shannon state entropy across that transition. Importantly, it is an exact difference. Therefore, for every closed directed cycle $C$,
\[
\boxed{
\sum_{e\in C}\Delta H_e=0.
}
\]
The identity follows by telescoping and has a direct methodological consequence: a state entropy difference can contribute to an open transition, but it cannot by itself generate a non-zero circulation around a closed cycle.

\section{Complex relational state and phase transport}

The same relational state is also represented by a normalized complex vector
\[
\psi(s)\in\mathbb C^m,
\qquad
\langle\psi|\psi\rangle=1.
\]
Whenever two consecutive states have non-zero overlap, define the normalized Pancharatnam link, following the standard geometric-phase construction \cite{Pancharatnam1956,Berry1984},
\[
G_{a\to b}
=
\frac{\langle\psi_a|\psi_b\rangle}
{|\langle\psi_a|\psi_b\rangle|}
\in U(1).
\]
An open link is gauge-covariant under local rephasing of the state representatives. The ordered product of links around a closed cycle is gauge-invariant and defines the discrete geometric holonomy
\[
\gamma_B(C)
=
\operatorname{Arg}\!\left(\prod_{e\in C}G_e\right).
\]

\section{Shannon--phase transition link}

To connect informational change and phase transport without identifying them prematurely, introduce a transition-associated non-exact information increment $\sigma_e$. The first microscopic candidate closure is obtained from forward/reverse transition asymmetry. Log-ratios of directed transition weights and their cycle affinities are standard objects in network nonequilibrium thermodynamics \cite{Schnakenberg1976}; the IDT-specific step is their typed coupling to the temporal phase link. For strictly positive admitted transition probabilities,
\[
\boxed{
\sigma_{a\to b}
=
\log_2\frac{P(b\mid a)}{P(a\mid b)}.
}
\]
This quantity is path-level rather than a state entropy. It obeys $\sigma_{b\to a}=-\sigma_{a\to b}$ and vanishes for a symmetric transition pair. The candidate temporal link is
\[
\boxed{
L_e
=
G_e\exp\!\left[i\kappa\left(\Delta H_e+\sigma_e\right)\right].
}
\]
The corresponding transition phase weight is
\[
q_e=\operatorname{Arg}L_e.
\]
It can therefore enter the atomic temporal transition measure used later in the construction,
\[
\mathcal K_T=\sum_e q_e\,\delta_{s_e}.
\]

For a closed cycle the exact Shannon contribution cancels. Defining the cycle affinity
\[
\mathcal A_C
=
\sum_{e\in C}\sigma_e
=
\log_2\prod_{e\in C}
\frac{P_{e,+}}{P_{e,-}},
\]
gives
\[
\boxed{
\operatorname{Arg}\!\left(\prod_{e\in C}L_e\right)
=
\gamma_B(C)
+
\kappa\mathcal A_C
\pmod{2\pi}.
}
\]
This is the first structural separation in the derivation. The geometric part supplies phase holonomy. The exact Shannon state potential supplies open-transition information change. Directed path affinity supplies a non-exact information channel that can survive cycle closure.

A useful control follows immediately. If a strictly positive relational weight $\pi_s$ satisfies pairwise detailed balance,
\[
\pi_aP(b\mid a)=\pi_bP(a\mid b),
\]
then
\[
\sigma_{a\to b}=\log_2\frac{\pi_b}{\pi_a}
\]
and the cycle affinity telescopes to zero. Non-zero $\mathcal A_C$ is therefore a precise obstruction to this detailed-balance representation on the declared cycle.

\section{Relational density, viscosity, and mobility}

Let the relational state carry strictly positive scalar fields $\rho_R(s)$ and $\eta_R(s)$. The minimal pair mobility used here is
\[
\boxed{
M_{ab}
=
\frac{\sqrt{\rho_R(a)\rho_R(b)}}
{\tfrac12[\eta_R(a)+\eta_R(b)]}.
}
\]
It is symmetric under $a\leftrightarrow b$. To keep pace and direction as separate typed quantities, introduce an antisymmetric dimensionless edge drive
\[
A_{ab}=-A_{ba}
\]
and define
\[
\boxed{
W_{a\to b}=M_{ab}e^{A_{ab}/2},
\qquad
W_{b\to a}=M_{ab}e^{-A_{ab}/2}.
}
\]
Then
\[
\sigma_{a\to b}
=
\log_2\frac{W_{a\to b}}{W_{b\to a}}
=
\frac{A_{ab}}{\ln2}.
\]
The symmetric mobility cancels from the directional ratio. Thus, in this minimal construction, relational density and viscosity control transition pace while the antisymmetric edge drive controls directional affinity.

A second structural control follows. If the edge drive is an exact state-potential difference,
\[
A_{ab}=V_R(b)-V_R(a),
\]
then its closed-cycle sum vanishes. Therefore a non-zero temporal circulation requires a non-exact edge drive or an enlarged state in which memory/history changes the relevant transition geometry.

The rate pair also admits a symmetric/antisymmetric decomposition,
\[
\mathfrak a_{ab}=W_{a\to b}+W_{b\to a},
\qquad
\mathfrak j_{ab}=W_{a\to b}-W_{b\to a}.
\]
The activity $\mathfrak a_{ab}$ is strictly positive and invariant under edge reversal, while the current $\mathfrak j_{ab}$ reverses sign. This decomposition will be used in the NOW layer to distinguish event existence from directional current.

The next derivational target is consequently sharper: identify whether memory-orbit state, phase holonomy, or another admitted relational structure generates the non-exact part of $A_{ab}$, while $\rho_R$ and $\eta_R$ regulate the pace of realization.

% ---- END monograph/chapters/01_information.tex ----

% ---- BEGIN monograph/chapters/03_kahler_information_geometry.tex ----
\chapter{Kähler Information Geometry}
The temporal construction begins on a state/information manifold \(\mathcal M\). Spacetime coordinates enter later at Einstein Closure.

\section{Relational forcing of the local geometry}
A positive context-sensitive internal elapsed measure has the local form
\[
 d\tau_{\rm int}=\phi(x,\lambda)\,d\lambda,
 \qquad \phi>0,
\]
so the scalar temporal pace is the Radon--Nikodym density of internal elapsed time with respect to the ordering parameter. In the current kinetic realization,
\[
\boxed{\phi=\frac{\mathfrak a}{\mathfrak a_\star}.}
\]

For a scalar informational functional \(\mathcal I\), the differential \(d\mathcal I\) is a covector. Coordinate-covariant local linear response therefore introduces a contravariant rank-two tensor \(G^{AB}\):
\[
V_I^A=-G^{AB}\partial_B\mathcal I.
\]
Positive informational descent makes the symmetric response positive semidefinite. On each active nondegenerate tangent sector it defines the inverse metric response.

A reversible phase response generated locally by \(\mathcal H_\alpha\) takes the form
\[
V_H^A=\Omega^{AB}\partial_B\mathcal H_\alpha,
\]
and preservation of \(\mathcal H_\alpha\) for arbitrary admitted phase covectors makes \(\Omega^{AB}\) antisymmetric. With the admitted unit \(U(1)\) phase generator \(J\),
\[
J^2=-I,
\qquad
\Omega^{AB}=J^A{}_{C}G^{CB},
\]
which yields the local exact-patch factorization
\[
\boxed{
\frac{dx}{d\lambda}
=-\operatorname{grad}_{g_K}\mathcal I
+J\operatorname{grad}_{g_K}\mathcal H_\alpha.
}
\]

\section{Shannon--Onsager realization of the symmetric response}
Formalism 01C fixes the stationary-reference informational scalar to the Kullback--Leibler divergence \cite{KullbackLeibler1951},
\[
\boxed{
\mathcal I_\pi[p]=D_{\rm KL}^{(2)}(p\|\pi).
}
\]
For an admitted stationary Markov kernel, data processing gives
\[
\mathcal I_\pi[pP]\le\mathcal I_\pi[p].
\]
The use of an informational potential in a dissipative response is consistent with the Onsager framework \cite{Onsager1931}; finite-state reversible Markov chains also admit rigorous entropy-gradient-flow formulations \cite{Maas2011}. The tensor below is the specific response adopted by the present formalism. In the continuous-time detailed-balance sector, let
\[
c_{ab}=\pi_aQ_{ab}=\pi_bQ_{ba},
\qquad
r_a=\frac{p_a}{\pi_a},
\]
and define the logarithmic mean
\[
\Lambda(x,y)=\frac{x-y}{\ln x-\ln y}.
\]
The exact symmetric response tensor is
\[
\boxed{
G^{(2)}_\pi(p)
=(\ln2)D^\top
\operatorname{diag}[c_{ab}\Lambda(r_a,r_b)]D,
}
\]
with
\[
\boxed{
\dot p=-G^{(2)}_\pi(p)\nabla_p\mathcal I_\pi.
}
\]
It obeys
\[
G^{(2)}_\pi\mathbf1=0,
\]
so the constant covector is the mass-conservation null. For a connected response graph, the metric is defined on the probability-simplex tangent quotient.

In the uniform zero-drive relational sector,
\[
\boxed{
G^{(2)}_u(u)=\frac{\ln2}{m}K_0,
\qquad
K_0=D^\top\operatorname{diag}(M_{ab})D.
}
\]
Thus the symmetric information geometry and the later Temporal Wave stiffness share the same relational-mobility Laplacian at uniform equilibrium.

\section{Global phase connection and holonomy}
The Shannon--Pancharatnam primitive assigns each admitted edge
\[
L_e=G_e\exp\!\left[i\kappa(\Delta H_e+\sigma_e)\right].
\]
Exact scalar differences obey \(\sum_C\Delta H_e=0\), while the closed-cycle temporal phase is
\[
\boxed{
\Phi_T(C)=\gamma_B(C)+\kappa\sum_{e\in C}\sigma_e\pmod{2\pi}.
}
\]
Admitted nonzero \(\Phi_T\) is carried globally by the temporal \(U(1)\) connection and its holonomy. The scalar \(\mathcal H_\alpha\) remains the local generator on an exact patch; the connection supplies the transition and cycle data needed to patch those local descriptions across the state bundle.

The associated two-form is
\[
\omega(X,Y)=g_K(JX,Y).
\]
The complex-projective/Fubini--Study sectors used by the Memory construction follow the standard quantum-state geometry \cite{ProvostVallee1980}, while the admitted geometric-phase sector is connected to the Pancharatnam--Berry construction \cite{Pancharatnam1956,Berry1984}. On compatible active sectors the resulting package
\[
\boxed{(\mathcal M,J,g_K,\omega)}
\]
is Kähler under the declared compatibility and closure conditions.

For an oriented two-chain \(\Sigma\), the signed Kähler-area diagnostic
\[
\mathcal A_K[\Sigma]=\int_\Sigma\omega
\]
obeys
\[
\mathcal A_K[-\Sigma]=-\mathcal A_K[\Sigma].
\]
This gives the formalism an orientation-sensitive geometric object. Temporal chirality is then tested downstream by the dynamics and holonomy gates.

\section{Visual phase portrait}
Figure~\ref{fig:kahler-flow-2d} gives a schematic 2D slice of the Kähler time flow. The nested contours represent levels of the informational functional; representative trajectories combine inward informational descent with the transverse phase response.
\begin{figure}[H]
  \centering
  \begin{tikzpicture}[scale=1.05]
    \draw[->] (-3.4,0) -- (3.4,0) node[right] {$q^1$};
    \draw[->] (0,-2.5) -- (0,2.5) node[above] {$q^2$};
    \draw[dashed] (0,0) ellipse (2.6 and 1.8);
    \draw[dashed] (0,0) ellipse (1.9 and 1.3);
    \draw[dashed] (0,0) ellipse (1.2 and 0.8);
    \draw[->,thick] (2.7,1.55)
      .. controls (1.75,2.05) and (0.65,1.55) .. (0.35,0.85)
      .. controls (0.05,0.20) and (0.65,-0.45) .. (1.00,-0.18)
      .. controls (1.25,0.03) and (0.80,0.45) .. (0.38,0.28);
    \draw[->,thick] (-2.7,-1.35)
      .. controls (-1.75,-1.95) and (-0.65,-1.45) .. (-0.35,-0.80)
      .. controls (-0.05,-0.15) and (-0.60,0.40) .. (-0.92,0.18)
      .. controls (-1.15,-0.02) and (-0.72,-0.40) .. (-0.32,-0.25);
    \draw[->,thick] (-2.45,1.50)
      .. controls (-1.20,1.80) and (0.30,1.20) .. (0.72,0.48)
      .. controls (1.02,-0.05) and (0.48,-0.58) .. (0.08,-0.38);
    \node[anchor=west] at (1.75,1.95) {informational contours};
    \node[anchor=west] at (1.55,-1.55) {descent + phase lift};
  \end{tikzpicture}
  \caption{Schematic two-dimensional Kähler time-flow slice with informational level contours and representative trajectories combining symmetric descent with the transverse phase response.}
  \label{fig:kahler-flow-2d}
\end{figure}

% ---- END monograph/chapters/03_kahler_information_geometry.tex ----
